% Worked Examples: Concept 3.9.5 - Nonlinear Effects
\documentclass[11pt,a4paper]{article}
\usepackage{worked-examples-template}

\title{\textbf{Worked Examples}\\[0.5em]
\large Concept 3.9.5: Nonlinear Effects\\[0.3em]
\normalsize \coursesubtitle{} -- \coursetitle}
\author{Assoc.\ Prof.\ William Robertson \and Dr Sean McGowan}
\date{\today}

\begin{document}

\maketitle
\tableofcontents
\newpage

\section{Concept 3.9.5: Nonlinear Systems}

\subsection{Example 1: Actuation Limits in a Voice Coil Servo System}

Physical actuators have fundamental limits on force, velocity, and acceleration that constrain achievable performance.

\paragraph{Problem:}
A voice coil servo system is modeled by:
\[
m\frac{\mathrm{d}^2x}{\mathrm{d}t^2} = F = k_I I
\]
where $m = 0.5$ kg is the mass, $k_I = 10$ N/A is the motor constant, and the drive has the following limits:
\begin{itemize}
\item Maximum current: $I_{\max} = 2$ A
\item Maximum voltage: $V_{\max} = 50$ V
\end{itemize}

(a) Calculate the maximum acceleration $a_{\max}$ and maximum velocity $v_{\max}$.

(b) Determine the minimum time to move the mass a distance $\ell = 0.5$ m, starting and ending at rest.

(c) Sketch the optimal velocity and acceleration profiles.

(d) Recalculate for $\ell = 2$ m and explain the difference in the solution.

\paragraph{Solution:}

\textbf{Step 1: Maximum acceleration and velocity}

Maximum acceleration is limited by the maximum current:
\[
a_{\max} = \frac{F_{\max}}{m} = \frac{k_I I_{\max}}{m} = \frac{10 \cdot 2}{0.5} = 40 \text{ m/s}^2
\]

Maximum velocity is limited by the back-EMF voltage constraint. For a voice coil, $V = k_I v$, so:
\[
v_{\max} = \frac{V_{\max}}{k_I} = \frac{50}{10} = 5 \text{ m/s}
\]

\textbf{Step 2: Determine if velocity saturation occurs}

The transition distance where velocity saturation begins is:
\[
\ell_{\text{crit}} = \frac{v_{\max}^2}{a_{\max}} = \frac{5^2}{40} = \frac{25}{40} = 0.625 \text{ m}
\]

For $\ell = 0.5$ m: Since $\ell < \ell_{\text{crit}}$, velocity saturation does \textbf{not} occur.

\textbf{Step 3: Minimum time for $\ell = 0.5$ m (no velocity saturation)}

When velocity is not saturated, the optimal strategy is:
\begin{itemize}
\item Accelerate at maximum from $t = 0$ to $t = t_1$ (reaching midpoint at $\ell/2$)
\item Decelerate at maximum from $t = t_1$ to $t = 2t_1$ (reaching $\ell$ at rest)
\end{itemize}

Using the formula for no velocity saturation:
\[
t = 2\sqrt{\frac{\ell}{a_{\max}}} = 2\sqrt{\frac{0.5}{40}} = 2\sqrt{0.0125} = 2 \cdot 0.1118 = 0.224 \text{ s}
\]

The peak velocity reached at the midpoint is:
\[
v_{\text{peak}} = a_{\max} \cdot \frac{t}{2} = 40 \cdot 0.112 = 4.48 \text{ m/s}
\]

Since $v_{\text{peak}} = 4.48 < v_{\max} = 5$ m/s, our assumption of no velocity saturation is confirmed.

\textbf{Step 4: Minimum time for $\ell = 2$ m (with velocity saturation)}

For $\ell = 2$ m: Since $\ell > \ell_{\text{crit}} = 0.625$ m, velocity \textbf{will saturate}.

The optimal strategy now has three phases:
\begin{enumerate}
\item Accelerate at $a_{\max}$ until reaching $v_{\max}$
\item Coast at constant $v_{\max}$ (zero acceleration)
\item Decelerate at $a_{\max}$ until rest
\end{enumerate}

Time to reach maximum velocity:
\[
t_1 = \frac{v_{\max}}{a_{\max}} = \frac{5}{40} = 0.125 \text{ s}
\]

Distance covered during acceleration (and deceleration):
\[
\ell_1 = \frac{v_{\max}^2}{2a_{\max}} = \frac{25}{80} = 0.3125 \text{ m}
\]

Distance at constant velocity:
\[
\ell_2 = \ell - 2\ell_1 = 2 - 2(0.3125) = 1.375 \text{ m}
\]

Time at constant velocity:
\[
t_2 = \frac{\ell_2}{v_{\max}} = \frac{1.375}{5} = 0.275 \text{ s}
\]

Total time:
\[
t_{\text{total}} = 2t_1 + t_2 = 2(0.125) + 0.275 = 0.525 \text{ s}
\]

Using the formula:
\[
t = \frac{\ell}{v_{\max}} + \frac{v_{\max}}{a_{\max}} = \frac{2}{5} + \frac{5}{40} = 0.4 + 0.125 = 0.525 \text{ s}
\]

\textbf{Step 5: Optimal profiles}

\textbf{For $\ell = 0.5$ m (no saturation):}
\begin{itemize}
\item Acceleration: $a(t) = +40$ m/s$^2$ for $0 < t < 0.112$ s, then $a(t) = -40$ m/s$^2$ for $0.112 < t < 0.224$ s
\item Velocity: Triangular profile, reaching peak of 4.48 m/s at $t = 0.112$ s
\end{itemize}

\textbf{For $\ell = 2$ m (with saturation):}
\begin{itemize}
\item Acceleration: $a(t) = +40$ m/s$^2$ for $0 < t < 0.125$ s, $a(t) = 0$ for $0.125 < t < 0.4$ s, $a(t) = -40$ m/s$^2$ for $0.4 < t < 0.525$ s
\item Velocity: Trapezoidal profile, reaching $v_{\max} = 5$ m/s and holding it for 0.275 s
\end{itemize}

\textbf{Key insights:}
\begin{itemize}
\item Actuation limits impose fundamental constraints on achievable performance
\item Short moves are acceleration-limited; long moves are velocity-limited
\item Optimal minimum-time control uses bang-bang acceleration (maximum or minimum)
\item These fundamental limits cannot be overcome by controller design alone
\end{itemize}

\subsection{Example 2: Friction-Induced Limit Cycles}

Nonlinear friction can cause oscillations even in well-designed linear control systems.

\paragraph{Problem:}
Consider a cart-pendulum balance system with state feedback. The cart has mass $M_c = 2$ kg and experiences Coulomb friction:
\[
F_{\text{friction}} = -\mu_f M_t g \operatorname{sgn}(v)
\]
where $\mu_f = 0.1$ is the coefficient of friction, $M_t = 2.5$ kg is the total mass (cart plus pendulum), $g = 9.81$ m/s$^2$, and $v$ is cart velocity.

The state feedback controller produces force $F_c$ to balance the pendulum. At equilibrium near $v = 0$, the controller attempts to apply small corrective forces.

(a) Calculate the magnitude of the friction force.

(b) Explain why friction can cause limit cycle oscillations.

(c) What is the minimum controller force needed to overcome static friction?

(d) Suggest methods to mitigate friction effects.

\paragraph{Solution:}

\textbf{Step 1: Friction force magnitude}

The Coulomb friction force magnitude is:
\[
|F_{\text{friction}}| = \mu_f M_t g = 0.1 \times 2.5 \times 9.81 = 2.45 \text{ N}
\]

This force opposes motion: $F_{\text{friction}} = -2.45 \operatorname{sgn}(v)$

\textbf{Step 2: Mechanism of limit cycles}

Coulomb friction creates a "dead zone" where small controller forces cannot produce motion:

\begin{enumerate}
\item \textbf{Near equilibrium:} The controller tries to apply a small force $|F_c| < 2.45$ N to correct a small position error
\item \textbf{Stiction:} If $|F_c| < |F_{\text{friction}}|$, the cart remains stuck ($v = 0$) despite the applied force
\item \textbf{Error accumulation:} While stuck, the integral action (if present) or accumulated state error increases the controller force
\item \textbf{Break-away:} Eventually $|F_c|$ exceeds 2.45 N, and the cart suddenly accelerates
\item \textbf{Overshoot:} Once moving, dynamic friction may be lower, causing the cart to overshoot the target
\item \textbf{Reversal:} The controller reverses force direction to correct the overshoot
\item \textbf{Cycle repeats:} The cart gets stuck again on the opposite side, creating oscillation
\end{enumerate}

\textbf{Step 3: Minimum force to overcome friction}

To initiate motion from rest:
\[
F_{\text{min}} > |F_{\text{friction}}| = 2.45 \text{ N}
\]

For the system to respond to small disturbances or reference changes, the controller must be capable of generating forces exceeding 2.45 N. However, this creates a trade-off:
\begin{itemize}
\item High controller gain $\Rightarrow$ can overcome friction but may cause violent motions and overshoot
\item Low controller gain $\Rightarrow$ smooth control but gets stuck in friction dead zone
\end{itemize}

\textbf{Step 4: Mitigation strategies}

Several approaches can reduce friction-induced oscillations:

\textbf{1. Friction compensation:}
Add a feedforward term to cancel the friction:
\[
F_{\text{total}} = F_c + \hat{F}_{\text{friction}} = F_c + 2.45 \operatorname{sgn}(v_{\text{ref}})
\]
Challenges: Requires accurate friction model and velocity measurement; sign function is discontinuous.

\textbf{2. Dead-band in control:}
If position error is below a threshold $\epsilon$, set controller output to zero:
\[
F_c = \begin{cases}
K(x_{\text{ref}} - x) & \text{if } |x_{\text{ref}} - x| > \epsilon \\
0 & \text{if } |x_{\text{ref}} - x| \leq \epsilon
\end{cases}
\]
Allows small steady-state error to avoid constant chattering.

\textbf{3. Dither signal:}
Add a small high-frequency oscillation to "shake" the system out of stiction:
\[
F_c = F_{\text{linear}} + A\sin(\omega_d t)
\]
The dither amplitude $A$ should be small enough not to disturb performance but large enough to reduce effective friction.

\textbf{4. Smooth friction model:}
Replace $\operatorname{sgn}(v)$ with a smooth approximation:
\[
F_{\text{friction}} \approx -\mu_f M_t g \tanh(\alpha v)
\]
where $\alpha$ controls the transition steepness. This makes the system more amenable to analysis and control design.

\textbf{Key insights:}
\begin{itemize}
\item Nonlinear friction effects cannot be eliminated by linear control design
\item Friction creates fundamental limits on achievable positioning accuracy
\item Practical control systems must account for actuator nonlinearities
\item Trade-offs exist between accuracy, smoothness, and robustness
\item Advanced techniques (adaptive control, sliding mode control) can better handle friction but add complexity
\end{itemize}
\end{document}
